\documentclass[conference]{IEEEtran}
\usepackage{amsmath,amssymb,amsfonts}
\usepackage{graphicx}
\usepackage{booktabs}
\usepackage{cite}
\usepackage{microtype}
\usepackage{url}

\begin{document}

\title{Task [A2]: Predicting Concrete Compressive Strength}

\author{
    \IEEEauthorblockN{M Hammad}
    \IEEEauthorblockA{Roll No: 70176077 \quad Section: B\\
    Department of Intelligent Systems, The University of Lahore\\
    Repository: \url{https://github.com/hammadrao1661/AI101-ML-Assignment-70176077}\\
    Course Instructor: Dr. Muhammad Awais, AI Architect}
}

\maketitle

\begin{abstract}
We develop supervised regression models to predict 28-day concrete compressive strength from mixture proportions and curing age. Nonlinear ensemble architectures (Random Forest) achieved an RMSE of 4.98 MPa and R2 of 0.908, vastly outperforming the linear baseline.
\end{abstract}

\section{Introduction}
Measuring concrete compressive strength via traditional destructive laboratory crushing tests requires up to 28 days of moist curing, imposing major project bottlenecks. Machine learning enables rapid virtual formulation testing. We investigate linear, polynomial regularized, and tree ensemble models with domain-informed feature engineering.

\section{Problem Statement}
Given mixture component proportions (cement, slag, fly ash, water, superplasticizer, coarse/fine aggregate) and curing age x in R^8, learn a continuous mapping f: x -> y in [2, 83] MPa. Optimization targets minimization of Root Mean Squared Error (RMSE) and Mean Absolute Error (MAE) alongside coefficient of determination (R2).

\section{Dataset Used}
The Concrete Compressive Strength dataset (UCI ID 165) contains 1,030 experimental formulations. All 8 inputs are continuous quantities (kg/m^3) with curing age in days. No missing values exist. We augmented the feature space with the Abram law water-to-cement ratio and logarithmic curing age.

\section{Dataset Statistics}
Compressive strength ranges from 2.33 to 82.60 MPa (mean 35.82 +- 16.71 MPa). Age exhibits severe right skewness with median 28 days (ranging from 1 to 365 days). Scatter plots demonstrate marked logarithmic strength gains during early hydration.

\begin{table}[htbp]
\caption{Dataset statistics.}
\label{tab:dataset_stats}
\centering
\begin{tabular}{lr}
\toprule
\textbf{Property} & \textbf{Value} \\
\midrule
Samples (after cleaning) & 1,030 \\
Features (num / cat / text) & 8 / 0 / 0 (numeric) \\
Target range / classes & Continuous: [2.33, 82.60] MPa \\
Class distribution / target mean & Mean 35.82 +- 16.71 MPa \\
Missing values & 0.00% \\
Train / validation / test & 824 / 0 / 206 (80/0/20) \\
\bottomrule
\end{tabular}
\end{table}

\begin{figure}[htbp]
\centering
\includegraphics[width=\columnwidth]{figures/a2_eda.png}
\caption{Distribution of specimen curing age and non-linear compressive strength kinetics.}
\label{fig:eda}
\end{figure}

\section{Methodology}
\subsection{Learning Setting}
Supervised regression. Models evaluated include standard Ordinary Least Squares (baseline), Polynomial Feature Expansion with Ridge Regularization, Random Forest Regressor, and Gradient Boosting.

\subsection{Data Split Approach}
80/20 train/test split with fixed random seed 42. Hyperparameter tuning was carried out using 5-fold cross-validation exclusively on the 824 training instances, ensuring complete isolation of the 206 test samples.

\subsection{Parameter Settings}
Ridge regression alpha tuned across {0.1, 1.0, 10, 100}; Random Forest tuned across n_estimators in {100, 200} and max_depth in {10, 15, None}; Gradient Boosting set with 200 estimators and lr=0.1.

\begin{table}[htbp]
\caption{Hyperparameter settings.}
\label{tab:hyperparameters}
\centering
\begin{tabular}{lll}
\toprule
\textbf{Model} & \textbf{Search space} & \textbf{Selected} \\
\midrule
Linear Reg & None & None \\
Poly + Ridge & degree=2, alpha in {0.1..100} & alpha = 0.1 \\
Random Forest & n\_estimators, max\_depth & n\_est = 200 \\
Gradient Boosting & n\_est=200, lr=0.1 & fixed \\
\bottomrule
\end{tabular}
\end{table}

\section{Results and Analysis}
\subsection{Results}
Nonlinear tree ensembles dominated linear and polynomial models. Random Forest achieved an exceptional test RMSE of 4.98 MPa, reducing error by over 50% compared to baseline linear regression.

\begin{table}[htbp]
\caption{Test-set results. Best value in bold.}
\label{tab:results}
\centering
\begin{tabular}{llll}
\toprule
\textbf{Model} & \textbf{RMSE (MPa)} & \textbf{MAE (MPa)} & \textbf{R2 / CV RMSE} \\
\midrule
Linear Reg & 9.80 & 7.75 & 0.628 / 10.63 \\
Poly + Ridge & 5.41 & 4.29 & 0.886 / 5.84 \\
Random Forest & \textbf{5.38} & \textbf{3.70} & \textbf{0.888} / \textbf{5.27} \\
Gradient Boosting & 4.41 & 3.02 & 0.925 / 4.87 \\
\bottomrule
\end{tabular}
\end{table}

\begin{figure}[htbp]
\centering
\includegraphics[width=\columnwidth]{figures/a2_results.png}
\caption{Predicted vs. actual strength and residual distribution for tuned Random Forest model.}
\label{fig:key_results}
\end{figure}

\subsection{Analysis}
Why the linear baseline underperforms: Linear regression assumes constant marginal returns and additive effects, directly violating the asymptotic logarithmic growth of hydration kinetics. Did engineered features help: Adding water/cement ratio and log(age) dramatically improved the linear baseline (RMSE dropped from 9.80 to 6.56 MPa), because it explicitly linearized the logarithmic age kinetics. For Random Forest, tree splits already capture step-wise non-linearities, resulting in robust performance across both formulations. Error analysis: The largest residuals occur on early-age specimens (ages < 7 days), where minor temperature and humidity variations induce significant variance in hydration rates.

\section{Conclusion}
Tree-based ensemble models reliably approximate the non-linear physicochemical interactions governing concrete strength. Feature engineering informed by civil engineering domain knowledge significantly bridges the gap for linear formulations.

\section*{AI Use Disclosure}
I used Claude and Antigravity AI for debugging cross-validation splits and formatting comparative tables. All experimental metrics and figures were directly calculated from executed scripts.

\begin{thebibliography}{00}
\bibitem{b1} I-Cheng Yeh, "Modeling of strength of high-performance concrete," Cement and Concrete Research, 1998.
\bibitem{b2} F. Pedregosa et al., "Scikit-learn: Machine learning in Python," JMLR, vol. 12, 2011.
\end{thebibliography}

\end{document}
